\subsubsection*{Conclusión de la comparativa}

\campo{¿Qué dispositivo presenta una mayor capacidad de procesamiento?}

Aquí la respuesta depende de a qué le llamemos \enquote{mayor}. Si hablamos de frecuencia de reloj y de arquitectura, el que gana es el Samsung Galaxy S26 Ultra: sus 2 núcleos Oryon V3 a 4,74 GHz son la CPU más rápida de las cinco y están hechos en 3 nm. Si hablamos de cantidad de RAM, gana el Pixel 10 Pro XL con 16 GB en todas sus versiones. Y si hablamos de memoria extra para la GPU, ahí está el A19 Pro del iPhone 17 Pro Max, con 6 núcleos y ray tracing por hardware. Lo que no tiene discusión es el Motorola Moto G47: 2,4 GHz y 6 nm, un chip pensado para que la batería dure, no para que la app vaya rápida.

\campo{¿Cuál tiene una pantalla de mayor tamaño?}

El Samsung Galaxy S26 Ultra y el iPhone 17 Pro Max empatan en 6,9 pulgadas, y en centímetros son prácticamente el mismo teléfono (163,6 mm frente a 163,4 mm de alto). Detrás vienen el OnePlus 13 con 6,82'', el Pixel 10 Pro XL con 6,8'' y el Motorola Moto G47 con 6,67''. Pero como ya decía, el dato que de verdad separa a estos teléfonos no es la diagonal: es que el Motorola tiene 264 ppi en un LCD de 8 bits y el OnePlus tiene 510 ppi en un OLED de 10 bits con 4.500 nits. Comida la pantalla más grande, te puedes encontrar con el peor panel de la lista.

\campo{¿Qué diferencias encuentras en la capacidad de las baterías?}

Veo tres diferencias que no son del mismo tipo. La de \textbf{capacidad}: gana el OnePlus 13 con 6.000 mAh (es de silicio-carbono), seguido del Pixel y el Motorola con 5.200 mAh, el iPhone con unos 5.088 y el Samsung con 5.000. La de \textbf{velocidad de carga}, que es donde se ve la brecha grande, de los 20 W del Motorola a los 80 W del OnePlus (100 W si compras su cargador aparte). Y una tercera que no me esperaba: \textbf{las declaraciones del fabricante no son comparables entre sí}. Samsung dice 31 horas de vídeo, Apple 39, Motorola 43 sin especificar el test y OnePlus directamente no publica ninguna cifra. Luego ese 43 del Motorola no significa que dure más que el Samsung.

\campo{¿Qué dispositivos ofrecen más opciones de conectividad?}

El más completo es el Samsung Galaxy S26 Ultra: 5G con mmWave, Wi-Fi 7, Bluetooth 6.0, NFC, UWB, USB 3.2, eSIM y satélite. El iPhone 17 Pro Max y el Pixel 10 Pro XL están casi a su nivel. El OnePlus 13 es el raro de la lista, porque trae mando IR y se queda sin UWB. Y el Motorola Moto G47 es el que más se aleja de todo: Wi-Fi 5 en vez de Wi-Fi 7, sin UWB, sin eSIM, sin satélite y USB 2.0.

\campo{¿Cómo podrían influir estas diferencias en la ejecución de una misma aplicación móvil?}

Esta es la que más me interesa, porque es la que se nota de verdad. Si ejecutamos la misma app en los cinco, lo que cambia no es que \enquote{vaya más o menos rápido}, sino que hay cosas enteras que existen en unos y directamente no existen en otros. Una app de pago con NFC funciona en cuatro de cinco. Una app que use el 5G para precargar vídeo tiene que tener un plan B para el Motorola. Una app con efectos gráficos irá bien en el Samsung y \enquote{a trompicones} en el Moto G47. Lo que aprendo con esta comparativa es que \textbf{más potencia y más espacio no van de la mano}, y que por eso hay que decidir a qué se le da prioridad.